\documentclass[a4paper,11pt]{article}
\usepackage[a4paper,margin=2cm]{geometry}
\usepackage[T1]{fontenc}
\usepackage[utf8]{inputenc}
\usepackage[ngerman,english]{babel}
\usepackage{lmodern,microtype,xcolor,parskip,enumitem}
\usepackage[most]{tcolorbox}
\usepackage{titlesec}
\usepackage[hidelinks]{hyperref}
\definecolor{HeaderBlueDark}{RGB}{70,110,170}
\definecolor{RowBlueLight}{RGB}{235,243,255}
\definecolor{RowBlueDark}{RGB}{220,233,250}
\definecolor{SoftGray}{RGB}{90,90,90}
\setlength{\parindent}{0pt}
\setlist[itemize]{leftmargin=1.4em,itemsep=0.35em,topsep=0.35em}
\linespread{1.06}
\titleformat{\section}[block]{\Large\bfseries\color{HeaderBlueDark}}{}{0pt}{}
\titlespacing{\section}{0pt}{1.3em}{0.8em}
\newtcolorbox{exbox}{enhanced,breakable,colback=RowBlueLight,colframe=RowBlueDark,boxrule=0.4pt,arc=1.5mm,left=1.2mm,right=1.2mm,top=1mm,bottom=1mm,title={Beispiele (Examples)},fonttitle=\bfseries,colbacktitle=RowBlueDark,coltitle=black}
\NewDocumentEnvironment{grammarentry}{mm+b}{\begin{tcolorbox}[enhanced,breakable,colback=white,colframe=HeaderBlueDark,boxrule=0.6pt,arc=1.5mm,left=1.5mm,right=1.5mm,top=1mm,bottom=1mm,colbacktitle=HeaderBlueDark,coltitle=white,fonttitle=\bfseries,title={#1\hfill\normalfont\footnotesize #2}]#3\end{tcolorbox}}{}
\newcommand{\GermanText}[1]{\textbf{Deutsch: }#1\par}
\newcommand{\EnglishText}[1]{{\small\color{SoftGray}\textbf{English: }#1\par}}
\newcommand{\ExampleItem}[2]{\item \textbf{#1}\\{\small\color{SoftGray}#2}}
\title{Das Fragewort „weswegen“ (The question word “weswegen”)}
\date{}
\begin{document}
\maketitle
\tableofcontents
\newpage
\section{Grundbedeutung und Wortart (Basic meaning and word class)}
\begin{grammarentry}{weswegen}{Interrogativadverb (kausal)}
\GermanText{\textbf{weswegen} ist ein kausales Interrogativadverb. Es fragt nach dem Grund oder Anlass: „aus welchem Grund?“.}
\EnglishText{\textbf{weswegen} is a causal interrogative adverb. It asks about a reason or occasion: “for what reason?”.}

\end{grammarentry}
\section{Satzstellung und Verwendung (Word order and use)}
\begin{grammarentry}{Direkte und indirekte Frage}{Fragewortstellung}
\GermanText{In einer direkten W-Frage steht das finite Verb gewöhnlich an zweiter Stelle nach dem Fragewort. In einer indirekten Frage steht es am Ende. Es ist weitgehend mit \textbf{warum} und \textbf{weshalb} austauschbar. In indirekten Fragen steht das finite Verb am Ende.}
\EnglishText{In a direct information question, the finite verb normally follows the question word in second position. In an indirect question, it goes at the end. It is largely interchangeable with \textbf{warum} and \textbf{weshalb}. In indirect questions, the finite verb goes at the end.}
\end{grammarentry}
\begin{exbox}\begin{itemize}
\ExampleItem{Weswegen brauchst du eine Bestätigung?}{Why do you need a confirmation?}
\ExampleItem{Er sagte nicht, weswegen er so spät gekommen war.}{He did not say why he had arrived so late.}
\end{itemize}\end{exbox}
\section{Zusammenfassung (Summary)}
\begin{grammarentry}{Merksatz}{kurz und wichtig}
\GermanText{Merksatz: \textbf{weswegen} fragt nach einem Grund oder Anlass.}
\EnglishText{Summary: \textbf{weswegen} asks about a reason or occasion.}
\end{grammarentry}
\end{document}
